\section{Exact interpolation in the summable coefficient space}
\label{sec:interpolation}

The finite columns of Section~\ref{sec:finite} approximate the interpolation
equations, but the auxiliary functions must use the genuine spectral
integrals.  We correct the finite lists in an unweighted \(\ell^1\) space,
holding their extra coordinates fixed.  The useful small quantity is not the
norm of the full transformed-jet operator: it is the norm of its outputs at
distant nodes.  This allows a finite-block inverse followed by a small tail
correction.

The elimination method adapts the finite-block and tail argument of
OpenAI~\citep[Section~3]{SharpPlanar2026}.  The operator, target data, and quantitative
estimates below are those of the present Gaussian problem.

Let
\[
 \mathcal E=\ell^1(\mathcal N;\R^2),\qquad
 \mathcal T=\mathcal N\setminus\mathcal F,
\]
where \(\mathcal F=\mathcal N\cap[1,168]\) is the finite node set from
Section~\ref{sec:finite}.  The two coordinates at each node are its \(c\)- and
\(d\)-coefficients.  On a pair of lists we use the sum of their norms.  A
subscript \(\mathcal F\) or \(\mathcal T\) restricts the corresponding input
or output; the paired finite block has 336 coordinates.  Every operator norm
is induced by these unweighted norms.

\begin{proposition}[Exact interpolation]
\label{prop:exact-interpolation}
For every real \(k\ge2.36\), there is a unique pair of positive-node lists in
\(\mathcal E\), with extra coordinates \(e_i=E_iq(k)\), whose full lists
\(\ell_i\) give
\begin{equation}
 \begin{aligned}
 H_1(n)&=T_k(n),& H_1'(n)&=T_k'(n),\\
 H_2(n)&=0,& H_2'(n)&=0
 \end{aligned}
 \qquad(n\in\mathcal N),
 \label{eq:exact-jets}
\end{equation}
where \(H_1=p_{\ell_1}+K_{p_{\ell_2}}\) and
\(H_2=p_{\ell_2}+K_{p_{\ell_1}}\).  Let \(\ell_i^{\mathrm{ref}}=X_iq(k)\),
extended by zero on \(\mathcal T\).  Then
\begin{equation}
 \sum_{i=1}^2\norm{\ell_i^{\mathrm{ref}}}_1<4.8,\qquad
 \sum_{i=1}^2\norm{\ell_i-\ell_i^{\mathrm{ref}}}_1<2\cdot10^{-8},\qquad
 \sum_{i=1}^2\norm{\ell_i}_1<5.
 \label{eq:exact-list-bounds}
\end{equation}
The functions \(H_1,H_2\) are entire and real on the real axis.  In
particular, \(H_1-T_k\) and \(H_2\) have zeros of order at least two at every
positive node, including the nodes that are not lattice radii.  Their
Gaussian-damped functions form the real radial Schwartz Fourier pair of
Lemma~\ref{lem:fourier-pair}.
\end{proposition}

The bounded inverse below gives uniqueness for summable coefficient lists
with fixed extras and jets prescribed on the larger periodic node set
\(\mathcal N\).  This differs from reconstruction using only the actual
triangular shells: Talebizadeh Sardari~\citep[Theorem~1.7]{Sardari2021} constructs infinitely many
linearly independent radial Schwartz functions whose values and radial
derivatives, and those of their Fourier transforms, vanish on all those
shells.

\begin{proof}
We first solve the exact operator equation on the full space.  We then bound
the size and residual of the finite reference lists and use the same block
equations to estimate their correction.

\paragraph{The coupled jet equation.}
Write \(a_i\in\mathcal E\) for the positive-node part of \(\ell_i\).  By
\eqref{eq:cardinal-jets} and \eqref{eq:jet-operator}, the desired equations are
\begin{equation}
 \mathcal L\binom{a_1}{a_2}
 :=\begin{pmatrix}I&-S\\-S&I\end{pmatrix}\binom{a_1}{a_2}
 =\binom{\mathcal JT_k+S_e e_2}{S_e e_1}.
 \label{eq:infinite-interpolation-system}
\end{equation}
Here \(S\) acts on positive-node inputs, while \(S_e\) acts on the three
extra coordinates.  Both are bounded by Lemma~\ref{lem:kernel-bounds}.  The
target belongs to \(\mathcal E\) as well.  Indeed,
\(T_k'(n)=-e^{-k(n-1)}\), and \eqref{eq:node-constant-bounds} gives
\begin{equation}
 \left|\frac{T_k(n)}{Q_n}\right|
 +\left|\frac{T_k'(n)-D_nT_k(n)}{Q_n}\right|
 \le\frac{1+(1+2.12)/k}{1.75}\,e^{-k(n-1)}.
 \label{eq:target-jet-bound}
\end{equation}
The right side is summable, uniformly after replacing \(k\) by 2.36.  We
will use the exact constant
\[
 c_*:=\frac{1+(1+2.12)/2.36}{1.75}=\frac{548}{413}<\frac43.
\]

\paragraph{The finite block.}
Let \(D=\widetilde S_{\mathcal F\mathcal F}\) and form the paired quadrature
block
\[
 \mathcal L_D=\begin{pmatrix}I&-D\\-D&I\end{pmatrix}.
\]
With \(\eta=.004^8\), submultiplicativity and
Proposition~\ref{prop:finite-certificate} give
\(\norm{D^{512}}=\norm{(D^{64})^8}<\eta\).  The finite geometric identities for the matrices
\(V_+,V_-\) in \eqref{eq:finite-columns} yield
\[
 R_+=(I-D)^{-1}=V_+(I-D^{512})^{-1},\qquad
 R_-=(I+D)^{-1}=V_-(I-D^{512})^{-1}.
\]
The remaining inverse is its norm-convergent Neumann series.  Diagonalizing
the two-list block into its sum and difference gives
\[
 \mathcal L_D^{-1}=\frac12
 \begin{pmatrix}
 R_++R_-&R_+-R_-\\
 R_+-R_-&R_++R_-
 \end{pmatrix}.
\]
Bounding the two output blocks in each column, and using the finite sum-norm
bound, gives
\begin{equation}
 \norm{\mathcal L_D^{-1}}
 \le\norm{R_+}+\norm{R_-}
 \le\frac{\norm{V_+}+\norm{V_-}}{1-\eta}
 <\frac{36}{1-\eta}<37.
 \label{eq:finite-quadrature-inverse}
\end{equation}

We now replace this quadrature block by the exact one.  Apply
Lemma~\ref{lem:quadrature-error} to orders zero and one with \(y=1\).
For each basis input, the errors in the two transformed moments are less~than~\(10^{-31}\).
In the jet projection their multipliers are bounded by
\(1/Q_n\) and \((1+|D_n|)/Q_n\), both less than 2.  In particular, every
entry of \(S_{\mathcal F\mathcal F}-D\), and every entry of
\(S_{\mathcal F e}-\widetilde S_{\mathcal F e}\), has absolute value less
than \(10^{-29}\).  There are 168 output coordinates.  The paired block is
off-diagonal, so its column norm has no extra factor of two.  Hence
\begin{equation}
 \norm{\mathcal L_{\mathcal F\mathcal F}-\mathcal L_D}
 <2\cdot10^{-27}.
 \label{eq:finite-operator-perturbation}
\end{equation}
Factoring \(\mathcal L_{\mathcal F\mathcal F}\) through \(\mathcal L_D\) and
using another Neumann series gives
\begin{equation}
 \norm{\mathcal L_{\mathcal F\mathcal F}^{-1}}
 <\frac{37}{1-37\cdot2\cdot10^{-27}}<40.
 \label{eq:finite-exact-inverse}
\end{equation}
This bound is for the exact finite block, not for the full inverse.

\paragraph{Solving after elimination of the finite block.}
The scalar bound \eqref{eq:scalar-error-budgets} gives
\[
 U_{168}(1,0)<1.1\cdot10^{-14}=:\delta.
\]
Because \eqref{eq:tail-bound} is uniform in the input node locations, it
bounds outputs on \(\mathcal T\) from either finite or infinite positive-node
inputs.  Again the two off-diagonal blocks introduce no extra factor in the
combined \(\ell^1\) norm.  We obtain
\begin{equation}
 \norm{\mathcal L_{\mathcal T\mathcal F}}<\delta,\qquad
 \norm{\mathcal L_{\mathcal T\mathcal T}-I}<\delta,\qquad
 \norm{\mathcal L_{\mathcal F\mathcal T}}<4000.
 \label{eq:tail-operator-blocks}
\end{equation}
The last bound need not be small.  We first solve the finite equation for
its finite unknowns in terms of the tail unknowns.  Substituting that solution
into the tail equation leaves the Schur complement
\[
 \mathcal Q=\mathcal L_{\mathcal T\mathcal T}
 -\mathcal L_{\mathcal T\mathcal F}
  \mathcal L_{\mathcal F\mathcal F}^{-1}
  \mathcal L_{\mathcal F\mathcal T}.
\]
It satisfies
\[
 \norm{\mathcal Q-I}
 <\delta(1+40\cdot4000)=1.760011\cdot10^{-9}<1,
\]
so its Neumann series converges.  More explicitly, for any right side
\(\psi=(\psi_{\mathcal F},\psi_{\mathcal T})\), the solution is reconstructed
by
\begin{align*}
 a_{\mathcal T}
 &=\mathcal Q^{-1}\bigl(\psi_{\mathcal T}
   -\mathcal L_{\mathcal T\mathcal F}
    \mathcal L_{\mathcal F\mathcal F}^{-1}\psi_{\mathcal F}\bigr),\\
 a_{\mathcal F}
 &=\mathcal L_{\mathcal F\mathcal F}^{-1}
   \bigl(\psi_{\mathcal F}-\mathcal L_{\mathcal F\mathcal T}a_{\mathcal T}\bigr).
\end{align*}
Every operator in these formulas is bounded.  Thus \(\mathcal L\) is a
boundedly invertible operator on the full pair space \(\mathcal E^2\), not
merely on a truncation.  Applied to \eqref{eq:infinite-interpolation-system},
this proves existence and uniqueness for the specified extras.
We henceforth write \(a=(a_1,a_2)\) for this solution.

\paragraph{Size of the reference lists.}
The same-column bounds in Proposition~\ref{prop:finite-certificate}, applied
to the nonnegative entries of \(q(k)\), give
\begin{align}
 \sum_i\norm{\ell_i^{\mathrm{ref}}}_1
 &\le\left(2.3+\frac{5.6}{k}\right)
       \left(1+e^{-2k}+e^{-3k}+e^{-6k}+e^{-8k}\right)\notag\\
 &\le\left(2.3+\frac{5.6}{2.36}\right)
       \left(1+e^{-4.72}+e^{-7.08}+e^{-14.16}+e^{-18.88}\right)
 <4.8.
 \label{eq:reference-norm}
\end{align}
The last comparison is the corresponding scalar bound in
\eqref{eq:scalar-error-budgets}.  The norm includes all extras, so
\(\norm{e_1}_1+\norm{e_2}_1<4.8\).  The constant coordinates in
\eqref{eq:extra-data} satisfy the sharper bound
\begin{equation}
 |C_1|+|C_2|
 =\left|-.0074+\frac{.0115}{k}\right|
  +\left|.0096-\frac{.017}{k}\right|<.02.
 \label{eq:extra-constant-bound}
\end{equation}
Indeed, for \(k\ge2.36\), the first expression inside absolute values is
negative and the second is positive; their absolute values are less than
\(.0074\) and \(.0096\), respectively.  The exact lists have the same
constants because their extras are fixed.

\paragraph{The three finite residual terms.}
Let \(a^{\mathrm{ref}}\) be the paired positive-node part of the reference
lists, and let \(r\) be the right side of
\eqref{eq:infinite-interpolation-system} minus \(\mathcal L a^{\mathrm{ref}}\).
Put \(\gamma_i=g_iq(k)\), with \(g_i\) as in
\eqref{eq:finite-columns}, and define
\[
 \Delta S_{\mathcal F}
 =\bigl(S_{\mathcal F\mathcal F}-D,\,
        S_{\mathcal F e}-\widetilde S_{\mathcal F e}\bigr).
\]
This map acts on the finite positive-node coordinates followed by the three
extras; the reference lists are supported on these coordinates.  The finite
geometric identities \((I-D)V_+=(I+D)V_-=I-D^{512}\) give the exact
decomposition
\begin{equation}\label{eq:finite-residual-decomposition}
 r_{\mathcal F}
 =\binom{\mathcal J_{\mathcal F}T_k-Yq(k)}{0}
  +\binom{\Delta S_{\mathcal F}\ell_2^{\mathrm{ref}}}
         {\Delta S_{\mathcal F}\ell_1^{\mathrm{ref}}}
  +\binom{D^{512}\gamma_1}{D^{512}\gamma_2}.
\end{equation}
The three terms are the omitted target data, the quadrature error on the
full reference lists, and the finite geometric-sum residual.

The vector \(Yq(k)\) reproduces the target jets at \(1,3,4,7,9\) and vanishes
at the other finite nodes.  Summing \eqref{eq:target-jet-bound} over all
integers from 12, a superset of the omitted nodes, and applying the scalar
bound in \eqref{eq:scalar-error-budgets}, gives
\begin{equation}
 \frac{1+(1+2.12)/2.36}{1.75}
 \sum_{n=12}^{\infty}e^{-2.36(n-1)}<8\cdot10^{-12}.
 \label{eq:omitted-target-jets}
\end{equation}
For the quadrature replacement, each basis column has finite output error
less than \(168\cdot10^{-29}\).  Multiplication by the combined reference
norm in \eqref{eq:reference-norm} therefore costs less than \(10^{-24}\),
including the extra inputs.

For the third term, we can bound \(\gamma_i\) without a further matrix
estimate.  The elementary
bound \(e^{-2.36}<2/21\), proved in the local calculation preceding
\eqref{eq:target-tail-168}, and \eqref{eq:target-jet-bound} give the first
five target jets total norm at most
\[
 c_*\left(1+(2/21)^2+(2/21)^3+(2/21)^6+(2/21)^8\right)<1.35.
\]
The exact operator on the extra inputs has norm less than 4000, and its finite
quadrature replacement differs by less than \(2\cdot10^{-27}\).  Since the
extras have combined norm less than 4.8,
\[
 \norm{\gamma_1}_1+\norm{\gamma_2}_1
 <1.35+(4000+2\cdot10^{-27})4.8<20000.
\]
The geometric-sum residual is consequently less than
\(20000\eta=1.31072\cdot10^{-15}<1.4\cdot10^{-15}\).  Adding the three
finite errors yields
\begin{equation}
 \norm{r_{\mathcal F}}_1<10^{-11}.
 \label{eq:finite-reference-residual}
\end{equation}

\paragraph{The tail residual.}
The positive-node reference lists vanish on \(\mathcal T\).  Their
transformed outputs there, including the extras, are bounded by
\(U_{168}(5,.02)\), using \eqref{eq:reference-norm} and
\eqref{eq:extra-constant-bound}.  For the target part, \(\mathcal T\) is
contained in the integers above 168, so \eqref{eq:target-jet-bound} gives
\[
 \sum_{n\in\mathcal T}
 \left(\left|\frac{T_k(n)}{Q_n}\right|
 +\left|\frac{T_k'(n)-D_nT_k(n)}{Q_n}\right|\right)
 \le c_*\frac{e^{-2.36\cdot168}}{1-e^{-2.36}}
 <10^{-170}.
\]
The last inequality is the local target-tail calculation
\eqref{eq:target-tail-168}.  The scalar bound
\eqref{eq:scalar-error-budgets} retains the separate strict margin
\(U_{168}(5,.02)<8.3\cdot10^{-14}-10^{-170}\).  Therefore
\begin{equation}
 \norm{r_{\mathcal T}}_1
 \le U_{168}(5,.02)+10^{-170}<8.3\cdot10^{-14}.
 \label{eq:tail-reference-residual}
\end{equation}

\paragraph{Correction size and exact jets.}
Put \(u_{\mathcal F}=\norm{(a-a^{\mathrm{ref}})_{\mathcal F}}_1\) and
\(u_{\mathcal T}=\norm{(a-a^{\mathrm{ref}})_{\mathcal T}}_1\).  Take the finite
and tail blocks of \(\mathcal L(a-a^{\mathrm{ref}})=r\).  Using
\eqref{eq:finite-exact-inverse}, \eqref{eq:tail-operator-blocks},
\eqref{eq:finite-reference-residual}, and \eqref{eq:tail-reference-residual}
gives
\begin{equation}
 \begin{aligned}
 u_{\mathcal F}&\le40\bigl(10^{-11}+4000u_{\mathcal T}\bigr),\\
 u_{\mathcal T}&\le8.3\cdot10^{-14}
       +1.1\cdot10^{-14}(u_{\mathcal F}+u_{\mathcal T}).
 \end{aligned}
 \label{eq:correction-system}
\end{equation}
Substituting the first inequality into the second gives
\[
 u_{\mathcal T}
 \le\frac{8.3\cdot10^{-14}+4.4\cdot10^{-24}}
          {1-1.760011\cdot10^{-9}}
 <8.31\cdot10^{-14},\qquad
 u_{\mathcal F}<1.37\cdot10^{-8}.
\]
Their sum is less than \(2\cdot10^{-8}\).  The extras did not change, so this
is the full-list correction in \eqref{eq:exact-list-bounds}.  Adding
\eqref{eq:reference-norm} gives the combined exact norm less than 5.

Finally, \eqref{eq:infinite-interpolation-system} is exactly
\(\mathcal JH_1=\mathcal JT_k\) and \(\mathcal JH_2=0\).  The definition of
\(\mathcal J\) recovers both values and first derivatives at every node of
\(\mathcal N\), proving \eqref{eq:exact-jets}.  The functions \(p_{\ell_i}\)
and \(K_{p_{\ell_i}}\) are entire and real on the real axis by
Section~\ref{sec:fourier}; \(T_k\) is entire as well.  The jet identities are
therefore genuine double zeros.  Lemma~\ref{lem:fourier-pair} supplies the
Schwartz Fourier pair.
\end{proof}
